\begin{example}
\label{example-locally-nilpotent-not-nilpotent}
Let $R = k[x_n | n \in \mathbf{N}]$ be the polynomial ring in infinitely
many variables over a field $k$. Let $I$ be the ideal generated by
the elements $x_n^n$ for $n \in \mathbf{N}$ and $S = R/I$. Then the ideal
$J \subset S$ generated by the images of $x_n$, $n \in \mathbf{N}$
is locally nilpotent, but not nilpotent. The convention here is
$\mathbf N=\{1,2,3,\ldots\}$, so no relation with exponent zero occurs.
An element $y\in J$ has a finite expression
$y=\sum_{i=1}^t b_i\overline{x}_{n_i}$ with $b_i\in S$ and $n_i\geq1$.
Put $N=1+\sum_{i=1}^t(n_i-1)$. In the full multinomial expansion
of $y^N$, every term has an exponent at least $n_i$ on one factor
$\overline{x}_{n_i}$, and hence vanishes. The empty sum gives $y=0$.
Thus $J$ is locally nilpotent.
For each $n\geq1$ define a $k$-algebra map
$$
R\longrightarrow k[t]/(t^{n+1}),\qquad
x_{n+1}\longmapsto t,\qquad x_j\longmapsto0\quad(j\ne n+1).
$$
Every generator $x_j^j$ of $I$ maps to zero, so this map factors
through $S$. The element $\overline{x}_{n+1}^n\in J^n$ maps to
$t^n\ne0$, since $1,t,\ldots,t^n$ form a $k$-basis of the target.
Consequently $J^n\ne0$ for every $n\geq1$ and $J$ is not nilpotent.
\end{example}